\documentclass[10pt]{article}

\usepackage[a4paper,top=13mm,bottom=13mm,left=16mm,right=16mm]{geometry}
\usepackage[T1]{fontenc}
\usepackage[utf8]{inputenc}
\usepackage[spanish,es-nodecimaldot]{babel}
\usepackage{lmodern}
\usepackage{microtype}
\usepackage{xcolor}
\usepackage{tabularx}
\usepackage{array}
\usepackage{enumitem}
\usepackage{ragged2e}

\definecolor{burdeos}{HTML}{922B38}
\definecolor{burdeososcuro}{HTML}{651C27}
\definecolor{tinta}{HTML}{252A31}
\definecolor{gris}{HTML}{626C78}
\definecolor{papel}{HTML}{F7F5F2}
\definecolor{linea}{HTML}{D7CFCA}
\definecolor{verde}{HTML}{176B5B}

\color{tinta}
\setlength{\parindent}{0pt}
\setlength{\parskip}{2.4pt}
\setlist[itemize]{leftmargin=4.5mm,itemsep=0.7pt,topsep=1pt}
\setlist[enumerate]{leftmargin=5.5mm,itemsep=0.8pt,topsep=1pt}
\renewcommand{\familydefault}{\sfdefault}

\newcommand{\sectiontitle}[1]{%
  \vspace{1.2mm}%
  {\color{burdeos}\fontsize{11.5}{13}\selectfont\bfseries #1}\par
  \vspace{0.4mm}{\color{burdeos}\rule{\linewidth}{0.75pt}}\vspace{0.5mm}
}

\newcommand{\phase}[2]{%
  \par\noindent\fcolorbox{linea}{papel}{%
    \begin{minipage}[t]{0.952\linewidth}
      \vspace{0.55mm}\small
      {\color{burdeososcuro}\bfseries #1}\quad #2
      \vspace{0.55mm}
    \end{minipage}}\par
}

\begin{document}
\thispagestyle{empty}
\noindent\begin{minipage}[t][\textheight][t]{\textwidth}

\begin{minipage}[c]{0.68\textwidth}
  {\color{burdeososcuro}\fontsize{18.5}{20.5}\selectfont\bfseries
  Horarios universitarios\par CP-SAT frente a Simulated Annealing\par}
  \vspace{0.7mm}
  {\fontsize{11.5}{13}\selectfont\bfseries
  Revisi\'on cr\'itica y evaluaci\'on reproducible para CB-CTT\par}
\end{minipage}\hfill
\begin{minipage}[c]{0.285\textwidth}
  \raggedleft\small\color{gris}
  Metodolog\'ia y T\'ecnicas de Investigaci\'on\par
  M\'aster en Aprendizaje Autom\'atico y Datos Masivos\par
  UPM | Pr\'actica 1 | 2026--2027
\end{minipage}

\vspace{1.8mm}
{\color{burdeos}\rule{\textwidth}{1.5pt}}

\sectiontitle{Motivaci\'on y delimitaci\'on}
\small\justifying
La elaboraci\'on de horarios universitarios exige asignar clases a periodos y aulas sin conflictos y, al mismo tiempo, equilibrar criterios de calidad como capacidad, dispersi\'on semanal, compacidad curricular y estabilidad de aula. Se trata de un problema NP-hard con inter\'es tanto para la investigaci\'on en optimizaci\'on combinatoria como para la IA aplicada a planificaci\'on y selecci\'on de algoritmos.

El trabajo se centrar\'a en \textbf{Curriculum-Based Course Timetabling (CB-CTT)}, formulaci\'on UD2 de la competici\'on ITC-2007. Sus 21 instancias p\'ublicas, restricciones y funci\'on de coste permiten una evaluaci\'on com\'un y reproducible. ITC-2019, m\'as pr\'oximo al timetabling real por incorporar asignaci\'on de estudiantes y restricciones m\'as ricas, se discutir\'a como referencia y l\'imite de validez externa, pero no se implementar\'a en esta entrega.

\sectiontitle{Objetivo y preguntas de investigaci\'on}
El objetivo es \textbf{sintetizar el estado actual de los m\'etodos para CB-CTT y comparar, bajo condiciones homog\'eneas, un modelo CP-SAT y un algoritmo de Simulated Annealing (SA)}. El estudio responder\'a a tres preguntas:
\begin{enumerate}
  \item \textbf{RQ1.} \textquestiondown Qu\'e papel desempe\~nan hoy los m\'etodos exactos, metaheur\'isticos, h\'ibridos y asistidos por aprendizaje en university timetabling?
  \item \textbf{RQ2.} Con el mismo presupuesto computacional, \textquestiondown c\'omo difieren CP-SAT y SA en factibilidad, calidad, estabilidad y comportamiento a lo largo del tiempo?
  \item \textbf{RQ3.} \textquestiondown Qu\'e caracter\'isticas de las instancias parecen favorecer a cada enfoque?
\end{enumerate}

\sectiontitle{Posicionamiento respecto al estado del arte}
Las revisiones recientes muestran que la frontera actual no enfrenta simplemente m\'etodos exactos contra heur\'isticas: predominan b\'usquedas locales avanzadas, ALNS, descomposici\'on, matheur\'isticas y esquemas h\'ibridos. En 2026, un SA con el vecindario \textit{LectureKick} obtuvo resultados competitivos con el SOTA de CB-CTT, mientras que otro trabajo combin\'o inicializaci\'on CP-SAT con VND y cadenas de Kempe. Por ello, la contribuci\'on propuesta no ser\'a un nuevo algoritmo de SOTA, sino una \textbf{comparaci\'on reproducible y cr\'itica de dos baselines complementarios}.

El aprendizaje autom\'atico se analizar\'a en su papel real dentro del dominio: selecci\'on de algoritmos, ajuste de par\'ametros, predicci\'on de dificultad y gu\'ia de la b\'usqueda. No se entrenar\'a un modelo con solo 21 instancias, pues esa muestra no permite una evaluaci\'on fiable ni una separaci\'on adecuada entre entrenamiento y prueba.

\sectiontitle{Metodolog\'ia}
\phase{Fase 1. Revisi\'on estructurada.}{B\'usqueda de trabajos publicados entre 2015 y septiembre de 2026 en Scopus, Web of Science e IEEE Xplore, complementada con rastreo de referencias. Se registrar\'an variante, datos, algoritmo, presupuesto, m\'etricas, validaci\'on estad\'istica y disponibilidad de c\'odigo.}

\vspace{1.1mm}
\phase{Fase 2. Implementaci\'on com\'un.}{Un mismo parser, validador, funci\'on de coste y formato de salida servir\'a a tres m\'etodos: construcci\'on greedy, CP-SAT con el objetivo oficial y SA con movimientos, swaps y reparaci\'on documentada.}

\vspace{1.1mm}
\phase{Fase 3. Protocolo experimental.}{Las configuraciones se fijar\'an antes de la evaluaci\'on; el ajuste usar\'a instancias artificiales o un conjunto reservado. Todos los m\'etodos compartir\'an hardware, n\'umero de n\'ucleos y l\'imites temporales. Se ejecutar\'an al menos diez semillas por instancia para los componentes estoc\'asticos.}

\vspace{1.1mm}
\phase{Fase 4. An\'alisis.}{Se medir\'an factibilidad, penalizaci\'on, gap frente al mejor valor conocido, tiempo hasta la primera soluci\'on factible, estabilidad y curva calidad--tiempo. La instancia ser\'a la unidad estad\'istica; se usar\'an pruebas pareadas, correcci\'on de Holm, tama\~nos del efecto e intervalos bootstrap.}

\sectiontitle{Contribuci\'on esperada, viabilidad y l\'imites}
Se espera obtener: (i) una taxonom\'ia actualizada de los enfoques para CB-CTT; (ii) una infraestructura experimental verificable; y (iii) evidencia sobre las fortalezas relativas de CP-SAT y SA. El proyecto es viable para tres personas distribuyendo revisi\'on bibliogr\'afica, modelo exacto e infraestructura, y metaheur\'istica y experimentaci\'on. C\'odigo, configuraciones, semillas y resultados crudos se conservar\'an en el repositorio.

Las principales amenazas a la validez son el reducido n\'umero de instancias, su procedencia de una sola universidad, la sensibilidad al hardware y a la implementaci\'on, y la distancia entre ITC-2007 y los requisitos actuales. Estas limitaciones se declarar\'an y se evitar\'an generalizaciones m\'as all\'a de la evidencia.

\sectiontitle{Referencias iniciales}
{\fontsize{7.25}{8.3}\selectfont\raggedright
[1] Di Gaspero, McCollum y Schaerf. \textit{ITC-2007: Curriculum-Based Course Timetabling}, 2007.\quad
[2] De la Rosa-Rivera et al. \textit{Expert Systems with Applications}, 174, 114694, 2021.\quad
[3] Ceschia, Di Gaspero y Schaerf. \textit{European Journal of Operational Research}, 308, 1--18, 2023.\par
[4] Abdipoor et al. \textit{Intelligent Systems with Applications}, 19, 200253, 2023.\quad
[5] Gu et al. \textit{Computation}, 13(1), 10, 2025.\quad
[6] Ceschia et al. \textit{Computers \& Operations Research}, 195, 107621, 2026.\quad
[7] Nguyen. \textit{ICAART 2026}, DOI: 10.5220/0014437700004052.}

\vfill
{\color{linea}\rule{\textwidth}{0.6pt}}\vspace{0.5mm}
{\color{gris}\fontsize{7.5}{8.5}\selectfont
Propuesta de trabajo | Equipo de tres estudiantes \hfill Entrega final: art\'iculo IEEE de 10 p\'aginas, p\'oster y fuentes LaTeX}

\end{minipage}
\end{document}
